\newcommand{\auditoriaid}{A06 -- Descripción de los pasos 00--07}
% Preambulo comun de las auditorias del informe E1 (revision_e1/)
\documentclass[11pt,a4paper]{article}
\usepackage[utf8]{inputenc}
\usepackage[T1]{fontenc}
\usepackage[spanish,es-noshorthands]{babel}
\usepackage{lmodern}
\usepackage{microtype}
% Sin patrones de guionado en espanol en esta instalacion (ver A10): se
% desactiva el guionado automatico para evitar cortes con reglas del ingles.
\hyphenpenalty=10000
\exhyphenpenalty=10000
\usepackage[a4paper,margin=2.2cm]{geometry}
\usepackage{booktabs,array,longtable}
\usepackage{graphicx,caption,subcaption}
\usepackage{xcolor}
\usepackage{enumitem}
\usepackage{fancyhdr}
\usepackage{amsmath}
\usepackage{tikz}
\usetikzlibrary{arrows.meta,positioning}
\usepackage{natbib}
\usepackage{xurl}
\usepackage[hidelinks]{hyperref}
\urlstyle{tt}
\newcommand{\archivo}[1]{\url{#1}}
\newcolumntype{L}[1]{>{\raggedright\arraybackslash}p{#1}}
\definecolor{alta}{RGB}{180,30,30}
\definecolor{media}{RGB}{200,120,0}
\definecolor{baja}{RGB}{60,110,160}
\definecolor{cajafondo}{RGB}{245,247,250}
\newcommand{\sevA}{\textcolor{alta}{\textbf{Alta}}}
\newcommand{\sevM}{\textcolor{media}{\textbf{Media}}}
\newcommand{\sevB}{\textcolor{baja}{\textbf{Baja}}}
\setlength{\parskip}{0.35em}
\setlength{\emergencystretch}{3em}
\setlength{\parindent}{0pt}
\pagestyle{fancy}
\fancyhf{}
\fancyhead[L]{\small Auditoría del Informe del Entregable~1 -- \auditoriaid}
\fancyhead[R]{\small \thepage}
\renewcommand{\headrulewidth}{0.4pt}
% Entorno de hallazgos: ID | Severidad | Ubicacion | Hallazgo y evidencia | Correccion
\newenvironment{hallazgos}{%
\small
\begin{longtable}{@{}L{1.15cm}L{1.25cm}L{1.9cm}L{6.0cm}L{\dimexpr\textwidth-10.3cm-10.6\tabcolsep\relax}@{}}
\toprule
\textbf{ID} & \textbf{Sev.} & \textbf{Ubicación} & \textbf{Hallazgo y evidencia} & \textbf{Corrección aplicada} \\
\midrule\endfirsthead
\toprule
\textbf{ID} & \textbf{Sev.} & \textbf{Ubicación} & \textbf{Hallazgo y evidencia} & \textbf{Corrección aplicada} \\
\midrule\endhead
\bottomrule\endfoot
}{\end{longtable}}
% Macros compartidas con el informe revisado (las secciones propuestas las usan)
% Macros compartidas entre el informe revisado y las auditorias.
% Cifras verificadas contra los archivos del repositorio (ver A01/A04/A07):
\newcommand{\nUniverso}{102}   % source_id CMIP6 con tos/Omon en ESGF (00b, model_availability_priority.csv)
\newcommand{\nSeed}{41}        % publican los cuatro experimentos (config/models_seed_cmip6.csv)
\newcommand{\nSel}{40}         % completos + QC PASS (models_inventory_final.csv)
\newcommand{\nComb}{160}       % combinaciones modelo-experimento PASS (qc_report.csv)
\newcommand{\nNoEnc}{44}       % config/models_no_encontrado_44.csv
\newcommand{\nParcial}{18}     % 102 - 40 - 44
% Fecha de la portada: fija (no \today), editable en un solo lugar.
\newcommand{\fechaentrega}{28 de septiembre de 2026}
% Estado de codigo que documenta el E1 (antes de los cambios del E2)
\newcommand{\commitEone}{\texttt{8ee6184}}

% En las auditorias el texto propuesto se muestra fuera del informe: las
% referencias cruzadas se imprimen con el nombre de la seccion destino.
\makeatletter
\def\etq#1#2{\expandafter\def\csname etq@#1\endcsname{#2}}
\etq{sec:resumen}{Resumen}
\etq{sec:alcance}{Alcance}
\etq{sec:marco}{Revisión científica}
\etq{sec:antecedentes}{Antecedentes}
\etq{sec:referencias}{Literatura sobre TOE}
\etq{tab:referencias}{bibliografía TOE}
\etq{sec:similitudes}{Niño~3.4 frente a Niño~1+2}
\etq{sec:toe_metodos}{Elección de los métodos de TOE}
\etq{tab:metodos_toe}{métodos TOE}
\etq{eq:szabo_signal}{ecuación señal M1}
\etq{eq:szabo_V}{ecuación V(t)}
\etq{eq:szabo_T}{ecuación T(t)}
\etq{eq:szabo_toe}{ecuación TOE M1}
\etq{eq:tod}{ecuación ToD}
\etq{sec:datos}{Datos}
\etq{fig:cmip6_contribution}{contribución institucional}
\etq{sec:tabla_modelos}{Modelos seleccionados}
\etq{tab:modelos}{modelos}
\etq{sec:desviacion}{Selección de modelos}
\etq{sec:dominios_periodos}{Dominios y periodos}
\etq{fig:region_nino}{regiones Niño}
\etq{sec:arquitectura}{Arquitectura del pipeline}
\etq{fig:flujograma}{Flujograma}
\etq{sec:scripts}{Descripción de los pasos}
\etq{sec:paso00b}{Paso 00b}
\etq{sec:paso01}{Paso 01}
\etq{sec:paso02}{Paso 02}
\etq{sec:paso04}{Paso 04}
\etq{sec:paso05}{Paso 05}
\etq{sec:paso06}{Paso 06}
\etq{sec:paso07}{Paso 07}
\etq{sec:reporte}{Reporte de estado}
\etq{sec:resultados}{Resultados}
\etq{fig:qc}{matriz de estado}
\etq{tab:cuentas}{conteo de modelos}
\etq{fig:mapas2d}{mapas 2D}
\etq{fig:series}{series}
\etq{fig:boxplot}{boxplots}
\etq{sec:roadmap}{Próximos pasos}
\etq{tab:roadmap}{próximos pasos}
\etq{sec:conclusiones}{Conclusiones}
\etq{sec:recomendaciones}{Recomendaciones}
\makeatother

\makeatletter
\AtBeginDocument{\renewcommand{\ref}[1]{\@ifundefined{etq@#1}{\textit{[#1]}}{\textit{\csname etq@#1\endcsname}}}}
\makeatother
% Texto propuesto: secciones degradadas un nivel y en caja
\newenvironment{propuesto}{%
  \par\medskip\noindent\textcolor{baja}{\rule{\textwidth}{0.8pt}}\par
  \noindent{\small\textbf{Inicio del texto propuesto para el informe revisado}}\par\medskip
  \begingroup
  \let\section\subsection\let\subsection\subsubsection\let\subsubsection\paragraph
}{%
  \endgroup
  \par\noindent{\small\textbf{Fin del texto propuesto}}\par
  \noindent\textcolor{baja}{\rule{\textwidth}{0.8pt}}\par\medskip}
% Recuadro de actualizacion posterior (decisiones del autor)
\newcommand{\actualizacion}[1]{\par\medskip\noindent\fcolorbox{media}{media!8}{\parbox{\dimexpr\linewidth-2\fboxsep-2\fboxrule\relax}{\small\textbf{Actualización del 29-09-2026 (decisiones del autor).} #1}}\par\medskip}

\begin{document}

\begin{center}
{\Large\bfseries Auditoría A06: Descripción de los pasos 00--07 y reporte de estado}\\[0.4em]
{\normalsize Informe del Entregable~1, Sección~6 (\texttt{informe\_e1\_smnv3.tex}, commit \texttt{330fd2e})}\\[0.2em]
{\small Revisión: \fechaentrega}
\end{center}

\section{Objeto}

Se audita la Sección~6, que describe cada script del \emph{pipeline}. Cada afirmación técnica se contrastó con el código del commit \texttt{8ee6184} (\archivo{revision_e1/trabajo/e1_snapshot/scripts/}). Las afirmaciones sobre el resultado (cobertura temporal, meses completos) se contrastaron además con los 160 archivos de \archivo{data/processed/masked/}.

\actualizacion{La Sección~6 del informe se redujo a un párrafo con las acciones para obtener los datos y el flujograma. El texto corregido paso a paso que resulta de esta auditoría pasó al anexo (sección ``Qué hace cada paso''), con dos ajustes: la ventana vigente (0°--360°, 30°S--30°N) en el Paso~03 y el techo del control de calidad en 45\,°C (antes 39\,°C). El techo se justifica porque, con la franja tropical completa, algunos modelos alcanzan de forma plausible entre 39 y 41\,°C en el mar Rojo (16°N, 40°E) y el golfo Pérsico (26°N, 54°E) hacia 2096--2100 en SSP3-7.0 y SSP5-8.5; se verificó en UKESM1-0-LL, CAS-ESM2-0 y ACCESS-CM2. A continuación se muestran ambos textos: primero el de la Sección~6 del informe y luego el del anexo.}

\section{Matriz de verificación}

\begin{center}\small
\begin{longtable}{@{}L{1.3cm}L{7.2cm}L{\dimexpr\textwidth-8.5cm-4\tabcolsep\relax}@{}}
\toprule
Paso & Afirmación del informe & Evidencia en el código (8ee6184) \\
\midrule\endhead
00 & Verifica \texttt{cdo}, \texttt{wget}, \texttt{python3}+\texttt{requests}; crea carpetas; no instala & \texttt{command -v}, \texttt{import requests}, \texttt{mkdir -p} (l.~17--34). Correcto. \\
00b & Nodo principal ORNL; universo de 102; \texttt{hist-1950} con respaldo en nodos alternativos; grilla más gruesa; detección de cambio & \texttt{ESGF\_SEARCH\_URL = esgf-node.ornl.gov/proxy/search}; \texttt{pick\_coarsest\_grid}. Correcto, salvo ``cuatro experimentos SSP'' (son tres SSP; A06-1). \\
00c & Busca modelos en \archivo{files_MD/*.md} & La carpeta es externa al repositorio (A06-2). \\
01 & Nodo LLNL; prefiere \texttt{r1i1p1f1}; idempotencia todo-o-nada; reintento de paginación & \texttt{esgf-node.llnl.gov/esg-search/search}; \texttt{if "r1i1p1f1" in common}; \texttt{retry\_full\_on\_truncate=2}. Correcto. \\
02a & 3 reintentos por réplica con pausa & \texttt{DOWNLOAD\_RETRIES=3}, \texttt{sleep 3}. Correcto. \\
02b & Nodos CEDA, DKRZ, IPSL, NSC, NCI & \texttt{ALT\_ESGF\_SEARCH\_URLS} tiene 8: esos 5 más ESGF-West, ORNL y DKRZ-Metagrid (A06-3). \\
02 & ``\texttt{run.sh} pregunta'' antes de (b) y (c) & La pregunta está en \texttt{02\_download\_all\_sources.sh}, solo con \texttt{-t 0} y modelos pendientes (A06-4). \\
02 & Tope de 3 intentos entre corridas & \texttt{MAX\_SEARCH\_ATTEMPTS=3}. Correcto. \\
04 & Filtro por \texttt{grid\_label}; \texttt{selvar}; \texttt{gencon}/\texttt{genbil}; pesos por experimento de respaldo; calendario; unidades por valor ($\geq$100); verificación de meses & Todas presentes (\texttt{fix\_calendar}, \texttt{fix\_units}, \texttt{months\_match}). Correcto. \\
05 & Máscara desde ERSSTv5; filtro $|v|>400$ & \texttt{VALUE\_LIMIT = 400.0}. Correcto. \\
06 & Rango $-2$ a 39\,°C; tolerancia 10\,\% en ``ssp245/ssp585''; \texttt{historical} $>1950$; sin caché & \texttt{MIN\_SST, MAX\_SST = -2.0, 39.0}; \texttt{LENGTH\_TOLERANCE = 0.9} para \emph{todos} los SSP, incluido \texttt{ssp370} (A06-5). \\
07 & \texttt{selected=True} si los cuatro pasan QC & \texttt{required\_experiments <= pass\_set}. Correcto; el informe no describe la corrección en esta sección (A06-6). \\
\bottomrule
\end{longtable}
\end{center}

\section{Verificación sobre los archivos procesados}

\begin{center}\small
\begin{tabular}{@{}lll@{}}
\toprule
Experimento & Periodo en 38--39 modelos & Excepciones \\
\midrule
\texttt{historical} & 1850--2014 (39) & IITM-ESM: 1900--2014 \\
\texttt{ssp245} & 2015--2100 (38) & CAMS-CSM1-0 e IITM-ESM: 2015--2099 \\
\texttt{ssp370} & 2015--2100 (38) & CAMS-CSM1-0: 2099; IITM-ESM: 2098 \\
\texttt{ssp585} & 2015--2100 (38) & CAMS-CSM1-0 e IITM-ESM: 2015--2099 \\
\bottomrule
\end{tabular}
\end{center}

Los conteos de meses coinciden exactamente con los años de cada archivo (1980, 1380, 1032, 1020 y 1008 meses), lo que confirma que no hay huecos.

\section{Hallazgos}

\begin{hallazgos}
A06-1 & \sevM & 00b & ``exige además los cuatro experimentos SSP configurados bajo una misma \texttt{grid\_label} \ldots''. Los SSP configurados son tres; lo que se exige bajo una misma grilla son los cuatro experimentos (histórico + 3 SSP), según \texttt{pick\_coarsest\_grid}. & ``los tres escenarios SSP \ldots\ variantes de grilla que cubren los cuatro experimentos''. \\
A06-2 & \sevM & 00c & \archivo{files_MD/*.md} se presenta como parte del proyecto, pero la carpeta no está en el repositorio, así que el paso no es reproducible desde un \texttt{git clone}. & Se describe como carpeta local externa. \\
A06-3 & \sevM & 02b & Lista 5 nodos alternativos; el código usa 8. & Se nombran los 8, agrupados por interfaz. \\
A06-4 & \sevM & 02b & ``Antes de entrar a este paso y al (c), \texttt{run.sh} pregunta'': la pregunta la hace \texttt{02\_download\_all\_sources.sh}, solo en sesión interactiva y con modelos pendientes. & Corregido. \\
A06-5 & \sevM & 06 & ``tolerancia del 10\,\% en \texttt{ssp245}/\texttt{ssp585}'' omite \texttt{ssp370}. & ``en los escenarios SSP''. \\
A06-6 & \sevM & 07 & La corrección del criterio de selección (\texttt{dc0bf9d}) aparece solo en el Resumen; la sección del Paso~07 no la menciona y no tenía etiqueta para referirla. & Se describe en el Paso~07, que recibe la etiqueta \texttt{sec:paso07}. \\
A06-7 & \sevM & 04 & El ejemplo de SSP incompleto nombra solo a CAMS-CSM1-0; IITM-ESM también publica menos años (histórico desde 1900; escenarios hasta 2098--2099). & Se nombran ambos: son los dos únicos casos y explican los conteos del QC. \\
A06-8 & \sevM & 06 & La regla ``\texttt{historical} debe llegar más allá de 1950'' es laxa por sí sola; el informe no dice qué garantiza la cobertura completa. & Frase que remite a la verificación mes a mes del Paso~04 y al resultado comprobado. \\
A06-9 & \sevB & 00 & ``Corrección de diseño: versiones previas de este documento asumían AWS CLI \ldots''. Es historia de versiones del documento, no del método. & Una frase sobre por qué se usa HTTP directo. \\
A06-10 & \sevB & 04 & Un solo párrafo de 380 palabras con siete subpasos numerados en línea, difícil de seguir. & Lista numerada, una acción por ítem. \\
A06-11 & \sevB & Intro Sec.~6 & ``La Figura~\ldots\ esquematiza el flujo \ldots\ descritos en la Sección~\ref{sec:scripts}'': la sección se cita a sí misma. & Autorreferencia eliminada. \\
A06-12 & \sevB & 02c & Paréntesis largo con guiones anidados dentro del ítem~(c). & Dos oraciones. \\
\end{hallazgos}

\section{Extensión}

Original: 1628 palabras. Propuesto: 479 
palabras. La lista numerada del Paso~04 y la nueva descripción del Paso~07 compensan los recortes de historia de versiones.

\section{Texto propuesto}

\begin{propuesto}
% ============================================================
\section{Obtención y procesamiento de los datos}
\label{sec:scripts}
% ============================================================

Para obtener los datos se consultó el índice de ESGF y se identificaron los 102 modelos que publican \texttt{tos}/\texttt{Omon}; se retuvieron los que publican los cuatro experimentos y, de cada uno, la variante de grilla más gruesa y un miembro de ensamble común a los cuatro. Los archivos se descargaron en cascada desde el nodo principal de ESGF, los nodos alternativos de la federación y, como último recurso, Copernicus CDS, y ERSSTv5 se descargó de NOAA PSL. Todos los modelos se regrillaron a la grilla de ERSSTv5, se homogeneizaron calendario y unidades (comprobando el valor numérico y no el atributo declarado), se enmascararon con la máscara océano-tierra de ERSSTv5 y se filtraron valores físicamente imposibles. Finalmente, un control de calidad verificó el rango físico de la TSM ($-2$ a 45\,°C; el límite superior considera que la franja tropical incluye mares marginales cálidos, como el mar Rojo y el golfo Pérsico, donde algunos modelos alcanzan de forma plausible entre 39 y 41\,°C hacia 2100 en SSP3-7.0 y SSP5-8.5) y la cobertura temporal de cada archivo, y solo se seleccionaron los modelos con los cuatro experimentos aprobados. La Figura~\ref{fig:flujograma} resume la secuencia. La descripción de cada paso y de sus controles se encuentra en el documento de anexos.

\begin{figure}[htbp]
\centering
\begin{tikzpicture}[
  scale=0.95, transform shape,
  box/.style={draw, rounded corners, fill=green!80!white!10, align=center, text width=4.2cm, minimum height=.8cm, font=\small},
  boxaux/.style={draw, rounded corners, dashed, fill=green!20!cyan!15, align=center, text width=2.6cm, minimum height=1.1cm, font=\small},
  edg/.style={-{Latex[length=3mm]}, thick},
  edgd/.style={-{Latex[length=3mm]}, dashed, gray!80!black},
  every node/.style={font=\small}
]
\node[box] (S00) at (0,0) {00: Entorno};
\node[box] (S00b) at (0,-1.5) {00b: Modelos candidatos};
\node[boxaux] (S00c) at (5.5,-1.5) {00c\\Verif. literatura};
\node[box] (S01) at (0,-3) {01: Catálogo ESGF};
\node[box] (S02) at (-2.6,-4.5) {02: Descarga CMIP6};
\node[box] (S03) at (3.9,-4.5) {03: Descarga ERSSTv5};
\node[box] (S02b) at (-2.6,-6) {02b: Nodos ESGF alt.};
\node[box] (S02c) at (-2.6,-7.5) {02c: Copernicus CDS};
\node[box] (S04) at (0.5,-9) {04: Grilla común};
\node[box] (S05) at (0.5,-10.5) {05: Máscara océano-tierra};
\node[box] (S06) at (0.5,-12) {06: Control de calidad};
\node[box] (S07) at (0.5,-13.5) {07: Inventario final};
\node[box] (S08) at (0.5,-15) {Registro de modelos};
\node[box] (S09) at (0.5,-16.5) {Gráficos, reporte y paquete};
\draw[edg] (S00) -- (S00b);
\draw[edgd] (S00b) -- (S00c);
\draw[edg] (S00b) -- (S01);
\draw[edg] (S01) -- (S02);
\draw[edg] (S02) -- (S02b);
\draw[edg] (S02b) -- (S02c);
\draw[edgd] (S02b) to[bend right=60] (S01);
\draw[edgd] (S02c) to[bend right=80] (S01);
\draw[edg] (S02.west) -- ++(-1.5,0) |- (S04.west);
\draw[edg] (S03) -- (S04);
\draw[edg] (S04) -- (S05);
\draw[edg] (S05) -- (S06);
\draw[edg] (S06) -- (S07);
\draw[edg] (S07) -- (S08);
\draw[edg] (S08) -- (S09);
\end{tikzpicture}
\caption[Flujograma del \emph{pipeline} del Entregable~1.]{Flujograma del \emph{pipeline}. Las flechas sólidas siguen la secuencia automática de \texttt{run.sh}: buscar los modelos (00b, 01), descargar (02, 03), homogeneizar (04, 05) y verificar (06, 07). Las punteadas indican el diagnóstico manual 00c y la actualización del catálogo cuando un modelo se obtiene de una fuente alternativa (02b, 02c). Las dos últimas cajas se ejecutan en cada corrida.}
\label{fig:flujograma}
\end{figure}


\medskip\noindent\textbf{Texto del anexo (detalle por paso):}

% ============================================================
\section{Qué hace cada paso (detalle técnico)}
\label{sec:pasos}
% ============================================================

Esta sección describe la lógica de cada script por su nombre y ruta, sin reproducir el código, que está disponible en el repositorio \archivo{smn_toe}. El flujograma de la Sección~6 del informe muestra el orden de los pasos; aquí se describe qué hace cada uno y qué controles incorpora.

\subsection{Paso 00 -- \texttt{scripts/00\_setup\_env.sh}}
Verifica las dependencias (\texttt{cdo}, \texttt{wget}, \texttt{python3} con el paquete \texttt{requests}) y crea la estructura de carpetas de datos y registros. No instala nada: si falta una herramienta, falla al inicio con un mensaje claro. La adquisición de datos se hace por HTTP directo contra ESGF y NOAA PSL, porque el repositorio público de CMIP6 en la nube de Amazon está en formato Zarr, incompatible con CDO, y no existe un equivalente para ERSSTv5.

\subsection{Paso 00b -- \texttt{scripts/00b\_build\_model\_list.py}}
Consulta el índice del nodo principal de ESGF (ORNL) y enumera todos los modelos CMIP6 que publican \texttt{tos}/\texttt{Omon}, que son 102. Para cada uno exige \texttt{historical} (o su equivalente \texttt{hist-1950} de HighResMIP), que busca también en nodos alternativos si el nodo principal no lo tiene indexado. Exige además los tres escenarios SSP configurados y retiene solo las variantes de grilla (\archivo{grid_label}) que cubren los cuatro experimentos; entre ellas elige la de mayor \archivo{nominal_resolution}, es decir, la más gruesa, convertida a kilómetros si el vocabulario la expresa en grados. Los modelos que califican quedan en \archivo{config/models_seed_cmip6.csv}. Para todos los modelos inspeccionados, el script escribe además \archivo{informe/model_availability_priority.csv}, con la disponibilidad 0/1 por experimento. Antes de repetir el barrido completo, que tarda más de una hora, hace una consulta rápida del número de modelos disponibles y solo rehace el barrido si ese número cambió.

\subsection{Paso 00c -- \texttt{scripts/00c\_check\_paper\_models.py} (diagnóstico manual)}
Busca los nombres de modelos CMIP6 citados en los resúmenes de la bibliografía del proyecto (una carpeta local de archivos Markdown, externa al repositorio) y los compara con la lista del Paso~00b. Los citados que no están en la lista se escriben en \archivo{config/models_missing_from_esgf.csv}, como candidatos para las fuentes alternativas. No forma parte de la secuencia automática de \texttt{run.sh}.

\subsection{Paso 01 -- \texttt{scripts/01\_query\_esgf\_catalog.py}}
Para cada modelo de la lista, consulta el nodo LLNL de ESGF y determina el miembro de ensamble (\archivo{variant_label}) común a los cuatro experimentos, prefiriendo \texttt{r1i1p1f1}. Escribe el catálogo \archivo{data/interim/models_catalog_status.csv} (estado por modelo, fuente, miembro y número acumulado de intentos de búsqueda) y las direcciones de descarga en \archivo{data/interim/esgf_file_urls.json}. Si el catálogo ya existe, no vuelve a consultar ESGF; para forzar una consulta nueva hay que borrar ambos archivos.

La búsqueda de archivos recorre el índice de ESGF por páginas hasta traer la lista completa. Un corte de red a mitad del recorrido puede truncar la lista sin mostrar error: en pruebas, la misma búsqueda de \texttt{EC-Earth3 historical} devolvió 107, 121 o 165 archivos según la corrida (165 es el total real). Por eso, si el recorrido se interrumpe, se repite completo desde el inicio, aquí y en el Paso~02b.

\subsection{Paso 02 -- descarga en cascada}
\archivo{scripts/02_download_all_sources.sh} encadena tres fuentes, cada una como respaldo automático de la anterior:
\begin{enumerate}[label=(\alph*)]
\item \archivo{02_download_cmip6_chunks.sh}: descarga HTTP directa con las direcciones del Paso~01. Cada dirección se reintenta tres veces, con una pausa breve, antes de pasar a la siguiente réplica. Es necesario para modelos que publican un archivo por año (hasta 165 archivos solo para \texttt{historical}), donde un corte transitorio en algún archivo es probable.
\item \archivo{02b_search_alt_esgf_nodes.py}: para los modelos que el nodo principal no resolvió, repite la búsqueda en ocho índices alternativos de la federación (CEDA, DKRZ, IPSL, NSC y NCI con la interfaz clásica, y ORNL, DKRZ y ESGF-West con la interfaz Metagrid), porque un modelo puede faltar en un índice por problemas de replicación y estar en otro.
\item \archivo{02c_download_copernicus_cds.py}: último recurso, solo para los modelos de \archivo{config/models_copernicus_available.csv}, la lista de modelos del conjunto \texttt{projections-cmip6} de Copernicus. Estar en la lista no garantiza que CDS tenga todos los experimentos; muchos pedidos de modelos poco replicados fallan con \texttt{RoocsValueError}. Requiere credenciales personales (\texttt{\textasciitilde/.cdsapirc}); sin ellas el paso se omite sin detener el resto.
\end{enumerate}
Antes de las fuentes (b) y (c), si hay modelos pendientes y la sesión es interactiva, el script pregunta si se desea intentarlas. Espera 15\,s y, sin respuesta, continúa como si fuera afirmativa, de modo que nunca bloquea una corrida desatendida. \texttt{run.sh} invoca este paso a través de \archivo{scripts/run_download_if_idle.sh}, que no lanza una segunda descarga si ya hay una en curso. Un modelo que agota (b) y (c) sin resolverse suma un intento en el catálogo; con tres intentos en corridas distintas deja de reintentarse hasta que se borre el catálogo.

\subsection{Paso 03 -- \texttt{scripts/03\_download\_ersstv5.sh}}
Descarga ERSSTv5 desde NOAA PSL una sola vez y la recorta con CDO a la ventana de descarga: toda la franja tropical, 0°--360° y 30°S--30°N (\texttt{config/domains.yaml}).

\subsection{Paso 04 -- \texttt{scripts/04\_process\_to\_common\_grid.sh}}
Para cada modelo y experimento:
\begin{enumerate}[label=(\arabic*)]
\item Usa solo los archivos crudos de la variante de grilla registrada en el catálogo. Es necesario porque una misma carpeta puede contener archivos de dos grillas del mismo periodo (ocurre con \texttt{CESM2}, con \texttt{gn} y \texttt{gr}), y al fusionarlos cada mes quedaría duplicado.
\item Descarta toda variable distinta de \texttt{tos} antes de regrillar, porque algunos modelos de malla curvilínea, como IPSL-CM6A-LR, incluyen variables auxiliares que hacen fallar a CDO.
\item Calcula una vez por modelo los pesos de interpolación hacia la grilla de ERSSTv5: conservativa (\texttt{gencon}) para mallas no estructuradas, como la de AWI-CM-1-1-MR, y bilineal (\texttt{genbil}) para el resto. Si los pesos del modelo fallan en un experimento, se recalculan pesos propios para ese experimento.
\item Fusiona en el tiempo los fragmentos ya regrillados y recorta al rango de años del experimento.
\item Asigna calendario \texttt{standard} si falta el atributo.
\item Convierte Kelvin a grados Celsius según el valor y no según el atributo: resta 273{,}15 solo si la media en Niño~3.4 es mayor o igual a 100. La regla es necesaria porque hay archivos que declaran \texttt{degC} pero contienen valores en Kelvin, como uno de GISS-E2-1-G (\texttt{ssp585}).
\item Verifica que el resultado tenga exactamente la misma lista de meses (año y mes, no solo el conteo) que los archivos crudos de la grilla correcta; si no coincide, descarta el resultado para reprocesarlo en la corrida siguiente.
\end{enumerate}
La comparación se hace contra lo que publican los archivos descargados y no contra un calendario ideal. Así, un modelo que publica permanentemente menos años de los pedidos se acepta tal cual: CAMS-CSM1-0 termina sus escenarios en 2099, e IITM-ESM empieza \texttt{historical} en 1900 y termina sus escenarios en 2098 o 2099.

\subsection{Paso 05 -- \texttt{scripts/05\_apply\_ocean\_mask.py}}
Construye una máscara océano-tierra derivada de ERSSTv5: un punto es océano si tiene al menos un valor válido en algún mes del registro. No usa la fracción de tierra propia de cada modelo (\texttt{sftlf}), de modo que la observación y todos los modelos comparten exactamente los mismos puntos válidos, condición necesaria para compararlos punto a punto. Se implementa en Python (\texttt{netCDF4}/\texttt{numpy}) porque la operación equivalente en CDO, aplicada sobre el valor de relleno, no propaga correctamente la máscara y marca como océano todo el dominio, incluida la tierra. Además aplica un filtro por valor físico: todo punto con $|\text{valor}| > 400$ pasa a \texttt{NaN}, independientemente de la máscara. El filtro es necesario porque FGOALS-f3-L contiene, en puntos que ERSSTv5 considera océano, un valor de relleno ($\sim$10$^{35}$) que el regrillado no reconoce como faltante; gracias a este filtro el modelo puede incorporarse al inventario.

\subsection{Paso 06 -- \texttt{scripts/06\_qc\_checks.py}}
Control de calidad automático de cada archivo procesado:
\begin{itemize}
\item \textbf{Rango físico}: la TSM debe estar entre $-2$ y 45\,°C. El límite superior considera que la franja tropical incluye mares marginales cálidos: en el mar Rojo y el golfo Pérsico algunos modelos alcanzan, de forma plausible, entre 39 y 41\,°C hacia 2100 en SSP3-7.0 y SSP5-8.5.
\item \textbf{Cobertura temporal}: en los escenarios SSP se admite hasta un 10\,\% menos de los meses esperados; en \texttt{historical} se exige que el registro llegue más allá de 1950, porque algunos modelos válidos no empiezan en 1850 (IITM-ESM empieza en 1900).
\end{itemize}
El resultado se escribe en \archivo{data/processed/qc_report.csv}. Las métricas se recalculan para todos los archivos en cada corrida, sin caché: un resultado guardado por nombre de archivo podría quedar desactualizado cuando un archivo se reprocesa. Como el control tarda menos de un minuto, recalcular siempre es la opción más segura. La exactitud de la cobertura la garantiza la verificación mes a mes del Paso~04; en la corrida final, los 40 \texttt{historical} terminan en 2014 y todos los archivos tienen sus meses completos.

\subsection{Paso 07 -- \texttt{scripts/07\_build\_inventory\_report.py}}
Combina el resultado del Paso~06 con la resolución de grilla (obtenida con \texttt{cdo griddes}) en \archivo{data/processed/models_inventory_final.csv}. Un modelo queda \texttt{selected=True} solo si aprueban el control de calidad los cuatro experimentos requeridos. Los cuatro se exigen explícitamente, en lugar de compararlos con los experimentos intentados, para que un modelo con descarga parcial no pueda quedar seleccionado aunque lo descargado apruebe.

\subsection{Reporte de estado -- \texttt{scripts/generate\_status\_report.py}}
Al final de cada corrida de \texttt{run.sh}, sin importar el rango de pasos, se escribe \archivo{logs/run_report_<fecha>.txt}, que clasifica cada modelo como descargado completo, con descarga parcial, procesado, descargado pero sin procesar, o pendiente de fuentes alternativas.

\end{propuesto}

\end{document}
